% BHU PYQ -- Manually transcribed & error-corrected
\documentclass[11pt,a4paper]{article}

\usepackage[a4paper,top=2.5cm,bottom=2.5cm,left=2.8cm,right=2.8cm,headheight=14pt]{geometry}
\usepackage{amsmath,amssymb}
\usepackage[T1]{fontenc}
\usepackage[utf8]{inputenc}
\usepackage{lmodern}
\usepackage{microtype}
\usepackage{fancyhdr}
\usepackage{enumitem}
\usepackage{hyperref}
\usepackage{lastpage}
\usepackage{parskip}

\hypersetup{colorlinks=true,urlcolor=black,linkcolor=black,
  pdftitle={CHB-101: Structure and Bonding & Organic Chemistry-I -- BHU PYQ},pdfauthor={Banaras Hindu University}}

\pagestyle{fancy}\fancyhf{}
\renewcommand{\headrulewidth}{0.4pt}\renewcommand{\footrulewidth}{0.4pt}
\fancyhead[L]{\small   \textit{Banaras Hindu University}}
\fancyhead[R]{\small   \textit{CHB-101: Structure & Bonding & Organic Chemistry-I}}
\fancyfoot[C]{\small Page  \thepage\ of \pageref{LastPage}}


\newlist{parts}{enumerate}{2}
\setlist[parts,1]{label=(\alph*),leftmargin=*,itemsep=5pt,topsep=3pt}
\setlist[parts,2]{label=(\roman*),leftmargin=*,itemsep=3pt,topsep=2pt}

\newcommand{\pts}[1]{\hfill   \textit{[#1]}}


\begin{document}


\begin{center}
  {\large\bfseries BANARAS HINDU UNIVERSITY}\\[2pt]
  {
\normalsize B.Sc. (Hons.) Semester I Examination, 2023-24}\\[6pt]
  \rule{0.8\linewidth}{0.5pt}\\[6pt]
  {\large\bfseries Paper No. CHB-101: Structure \& Bonding and Organic Chemistry-I}\\[4pt]
  {
\normalsize Time: 3 Hours\hfill Maximum Marks: 70}
\end{center}
\rule{\linewidth}{0.4pt}
\small



\noindent   \textit{Instructions: Answer five questions in total, including Question~1 from each section which is compulsory. Use separate answer books for Section~A and Section~B.}

\normalsize
\bigskip


\begin{center}
  {\large\bfseries SECTION A}\\[2pt]
  {
\normalsize (Structure and Bonding -- Marks: 35)}
\end{center}
\rule{0.4\linewidth}{0.2pt}
\smallskip




\noindent   \textbf{Question 1.} Answer all parts of the following: \pts{5 $ \times$ 2 = 10}

\begin{parts}
  \item What do the `$+$' and `$-$' signs in the boundary surface diagrams of orbitals indicate?
  \item What is the physical significance of the radial wave function and the radial distribution function?
  \item Mention the limitations of the radius ratio rules in brief.
  \item $ \text{NH}_3$ and $ \text{NF}_3$ are pyramidal in shape, but the bond angle around the nitrogen center is different in the two cases. Why?
  \item What inferences can be obtained from the knowledge of HOMO and LUMO in molecules?
\end{parts}

\medskip\rule{\linewidth}{0.2pt}




\noindent   \textbf{Question 2.}

\begin{parts}
  \item Draw the radial distribution function versus $r$ for $3 \text{s}$, $3 \text{p}$, and $3 \text{d}$ orbitals, clearly depicting the salient differences in all three cases. \pts{6}
  \item Discuss the band theory of solids in brief. Explain conductors, insulators, and semiconductors using band gaps. \pts{6}
\end{parts}

\medskip\rule{\linewidth}{0.2pt}




\noindent   \textbf{Question 3.}

\begin{parts}
  \item Write the expression for the Laplacian operator in Cartesian coordinates. How are Cartesian coordinates related to polar coordinates in the Schr\"odinger wave equation? \pts{6}
  \item Explain the defects/limitations of simple Molecular Orbital Theory by drawing the energy level diagram of $ \text{CO}$ using homonuclear diatomic concepts. Draw the correct MO diagram for $ \text{CO}$. \pts{6}
\end{parts}

\medskip\rule{\linewidth}{0.2pt}




\noindent   \textbf{Question 4.}

\begin{parts}
  \item Write the molecular orbital configuration and draw the molecular orbital energy level diagram of $ \text{NO}$. Comment on the change in bond order and magnetic behavior when $ \text{NO}$ is oxidized to $ \text{NO}^+$. \pts{6}
  \item Derive the limiting radius ratio ($r^+/r^-$) for tetrahedral ionic structures (coordination number 4). (Given: $\cos 30^\circ \approx 0.866$, $\sin 54^\circ44' \approx 0.8164$). \pts{6}
\end{parts}


\newpage

\begin{center}
  {\large\bfseries SECTION B}\\[2pt]
  {
\normalsize (Organic Chemistry-I -- Marks: 35)}
\end{center}
\rule{0.4\linewidth}{0.2pt}
\smallskip




\noindent   \textbf{Question 5.} Answer all parts of the following: \pts{5 $ \times$ 2 = 10}

\begin{parts}
  \item Explain why neopentyl bromide undergoes an elimination reaction even though it does not have a $\beta$-hydrogen.
  \item Write down the Newman projections of $n$-butane having the most and least torsional strain.
  \item Give the name and structure of the initiators used in anionic and cationic polymerization.
  \item Complete the following reaction showing the mechanism:
    \[  \text{CH}_3 \text{-CH}_2 \text{-CH}_2 \text{-CH}_2 \text{-OH} \xrightarrow{ \text{Conc. H}_2 \text{SO}_4} ? \]
  \item Why does a polar protic solvent decrease the rate of an $ \text{S}_{ \text{N}}2$ reaction?
\end{parts}

\medskip\rule{\linewidth}{0.2pt}




\noindent   \textbf{Question 6.}

\begin{parts}
  \item Complete the following reactions, giving a suitable mechanism: \pts{2 $ \times$ 4 = 8}
  
\begin{parts}
    \item $ \text{CH}_3 \text{-C}\equiv \text{CH} \xrightarrow{ \text{dil. H}_2 \text{SO}_4,  \text{HgSO}_4} ? \xrightarrow{ \text{i. Na, ii. CH}_3 \text{I}} ?$
    \item   \textit{cis}-$2$-Butene $\xrightarrow{ \text{HBr, H}_2 \text{O}_2} ?$
  \end{parts}
  \item Predict the $R/S$ or $E/Z$ nomenclature of the following molecules where applicable: \pts{4}
  
\begin{parts}
    \item Chiral center in $2$-bromobutane
    \item $ \text{Ph(CH}_3 \text{)C=C(Br)H}$
    \item $ \text{CH}_3 \text{(Cl)C=C(Br)I}$
    \item Chiral center in $1$-chloro-$2$-propanol
  \end{parts}
\end{parts}

\medskip\rule{\linewidth}{0.2pt}




\noindent   \textbf{Question 7.}

\begin{parts}
  \item How do you prepare the following compounds from phenylmagnesium iodide ($ \text{PhMgI}$)? Write down stepwise mechanisms: \pts{6}
  
\begin{parts}
    \item Benzoic acid
    \item Bromobenzene
    \item Phenol
  \end{parts}
  \item How do you distinguish primary and secondary alcohols by the Lucas test (substitution method)? \pts{3}
  \item How do you prepare a $1,3$-diketone from ethyl acetoacetate? \pts{3}
\end{parts}

\medskip\rule{\linewidth}{0.2pt}




\noindent   \textbf{Question 8.}

\begin{parts}
  \item How do you prepare the following compounds from diethyl malonate? \pts{6}
  
\begin{parts}
    \item Glutaric acid
    \item Barbituric acid
  \end{parts}
  \item Predict the major products of the following reactions: \pts{4}
  
\begin{parts}
    \item $ \text{CH}_3 \text{-CH}_2 \text{-CH}_2 \text{-Br} +  \text{EtSNa} \xrightarrow{ \text{EtOH}} ?$
    \item $ \text{CH}_2 \text{=CH-CH=CH}_2 +  \text{Ethylene} \xrightarrow{100^\circ \text{C}} ?$
  \end{parts}
  \item How do you prepare adipic acid from cyclohexene? \pts{2}
\end{parts}

\vfill
\rule{\linewidth}{0.4pt}

\begin{center}\small
  \href{https://bhu-exam.vercel.app/}{Click here to visit our website}\\[4pt]
  \href{https://me-aryan.vercel.app/}{Click here to visit our contributor}\\[4pt]
  \href{https://quashan-me.vercel.app/}{Click here to visit our contributor}
\end{center}

\end{document}
